\documentclass[10pt,a4paper]{amsart}
\usepackage[margin=19mm]{geometry}
\usepackage{amsmath,amssymb,mathtools}
\usepackage[scr=rsfs,cal=euler]{mathalfa}
\usepackage{xfrac}
\usepackage[unicode,hidelinks,bookmarks=false]{hyperref}
\newcommand{\ie}{i.e.\ }
\newcommand{\cf}{cf.\ }
\newcommand{\resp}{respectively\ }
\newcommand{\Subsection}{\subsection}
\providecommand{\nobreakdash}{\nobreak\hbox{-}}
\setlength{\emergencystretch}{2em}
\multlinegap0pt
\usepackage{amssymb}
\usepackage[shortlabels]{enumitem}
\newtheorem{proposition}{Proposition}
\title{Intrinsic Finite-Step Characterizations of Discrete General Helices}
\author{Dae Won Yoon, Chul Woo Lee, Jae won Lee}
\date{Source: arXiv:2609.00673v1 (CC BY 4.0)}
\begin{document}
\maketitle

\noindent\emph{Source and licence.} Intrinsic Finite-Step Characterizations of Discrete General Helices. Version arXiv:2609.00673v1 (CC BY 4.0), distributed by arXiv under the Creative Commons Attribution 4.0 International licence, http://creativecommons.org/licenses/by/4.0/, as recorded in the arXiv licence metadata for this exact version and shown on the abstract page https://arxiv.org/abs/2609.00673v1. Full licence notice: \texttt{licenses/arxiv-2609.00673-CC-BY-4.0.txt}.\par\smallskip

\noindent\textit{Editorial scope.} Counted unit: the article's proposition prop:smooth-consistency (source0.tex lines 20306--20362). The article's complete original proof of it is delivered with it (source0.tex lines 20365--20527, one \texttt{proof} environment of 163 lines). No mathematics is written, repaired or re-derived: the statement and the proof are the article's own lines, in the article's own notation, and no retained line is substituted or re-flowed.  No further source-local context is needed: every cross-reference of every retained line resolves inside the two delivered spans.  Nothing was omitted from any retained span, so the package carries no omission record.  No retained line contains a citation, so the wrapper carries no bibliography and no bibliography entry.  No retained line contains a figure environment, an image-inclusion command or drawing code, so no image is delivered and the package declares no asset dependency.  Editorial formatting only, in addition: an AMS \texttt{amsart} preamble that restates the article's own declarations and macros used by the retained lines, namely the environments \texttt{proposition} on separate counters, together with the standard packages those lines need.  The retained environments print in this document as Proposition 1 in order of appearance, whereas the article prints the same items with the numbering of its own full document.  The complete native source of the article as distributed in the arXiv e-print for this exact version is bundled unchanged as \texttt{sources/209049/source0.tex}.
\begin{proposition}[Second-order smooth consistency]\label{prop:smooth-consistency}
Let \(\gamma\in C^6(I,\mathbb R^3)\) be a unit-speed space curve.
Assume that, on a compact subinterval \(J\Subset I\),
\[
\kappa(s)\ge \kappa_0>0,\qquad |\tau(s)|\ge\tau_0>0.
\]
Sample it uniformly by
\[
P_i=\gamma(s_i),\qquad s_i=s_0+ih.
\]
For the quotient indexed by $i$, call $C_i=[s_{i-1},s_{i+2}]$ its associated four-point cell. All quotient estimates below are uniform whenever $C_i\subset J$; estimates for the defect $D_i^h$, which involves two consecutive quotients, are uniform whenever $C_i\cup C_{i+1}=[s_{i-1},s_{i+3}]\subset J$.
Let \(\theta_i^h\) and \(\phi_i^h\) be the turning and signed torsion
angles of the chord polygon and set
\[
x_i^h=\tan\frac{\theta_i^h}{2}.
\]
There exists $h_0=h_0(\kappa_0,\tau_0,J)>0$ such that, for $0<h<h_0$ and all cells under consideration,
\[
0<\theta_i^h<\pi,\qquad \phi_i^h\notin\{0,\pi\}.
\]
Hence the sampled polygon lies in the generic Frenet branch and all quotients below are well defined.
With $m_i=s_i+h/2$, the midpoint convention compatible with the chord tangents gives the
asymptotic expansions
\[
x_i^h=\frac h2\kappa(s_i)+O(h^3),
\]
and
\[
\phi_i^h=h\tau\left(s_i+\frac h2\right)+O(h^3).
\]
Consequently,
\[
\frac{x_i^h+x_{i+1}^h\cos\phi_i^h}{\sin\phi_i^h}
=
\frac{\kappa}{\tau}\left(s_i+\frac h2\right)+O(h^2).
\]
If
\[
D_i^h:=
\frac{x_{i+1}^h+x_{i+2}^h\cos\phi_{i+1}^h}{\sin\phi_{i+1}^h}
-
\frac{x_i^h+x_{i+1}^h\cos\phi_i^h}{\sin\phi_i^h},
\]
then, uniformly on such cells,
\[
\boxed{
\frac{D_i^h}{h}
=
\left(\frac{\kappa}{\tau}\right)'(m_i)+O(h)
},
\qquad m_i=s_i+\frac h2.
\]
Thus the finite-step conservation law approaches $(\kappa/\tau)'=0$, equivalently \(\tau/\kappa=\mathrm{constant}\) on intervals where \(\kappa\tau\neq0\). The scalar quotient itself is second-order accurate, whereas its normalized finite difference is first-order accurate. If the underlying smooth curve is a general helix, then
\[
D_i^h=O(h^3).
\]
\end{proposition}

\begin{proof}
Put \(m_i=s_i+h/2\). Symmetric Taylor expansion of the chord about \(m_i\) gives
\[
\gamma(m_i+h/2)-\gamma(m_i-h/2)
=hT(m_i)+\frac{h^3}{24}T''(m_i)+O(h^5).
\]
Since
\[
T''=-\kappa^2T+\kappa'N+\kappa\tau B,
\]
normalization yields
\begin{equation}
T_i^h
=
T(m_i)+\frac{h^2}{24}
\bigl(\kappa'N+\kappa\tau B\bigr)(m_i)+O(h^4).
\label{eq:chord-expansion}
\end{equation}
The neighboring chord tangents are centered at \(s_i-h/2\) and
\(s_i+h/2\). Expanding them symmetrically about \(s_i\) gives
\[
T_i^h-T_{i-1}^h=h\kappa(s_i)N(s_i)+O(h^3).
\]
Because the two vectors are unit,
\[
\|T_i^h-T_{i-1}^h\|
=2\sin\frac{\theta_i^h}{2},
\]
and therefore
\begin{equation}
x_i^h=\tan\frac{\theta_i^h}{2}
=\frac h2\kappa(s_i)+O(h^3).
\label{eq:x-asymptotic}
\end{equation}
Keeping one further term in the same symmetric Taylor calculation gives
\[
x_i^h=\frac h2\kappa(s_i)+h^3A_x(s_i)+O(h^4)
\]
for a smooth coefficient $A_x$. Its explicit expression is not needed below; only the smooth dependence of this coefficient on the cell is used. The $C^6$ hypothesis supplies uniform control of the remainder in this refinement and in the $O(h^4)$ torsion-angle expansion below.

We next compute the binormal more precisely.  Expanding the two adjacent
chord tangents about \(s_i\), taking their cross product, and normalizing
gives
\begin{equation}
B_i^h
=
B(s_i)+h^2C_B(s_i)+O(h^4),
\label{eq:B-asymptotic}
\end{equation}
where
\begin{equation}
C_B
=
-\frac{\kappa\tau}{6}T
-\frac{2\kappa'\tau+\kappa\tau'}{12\kappa}N.
\label{eq:CB-explicit}
\end{equation}
Notice in particular that \(C_B\perp B\), as required by the unit-length
constraint to this order.

To extract the signed torsion angle, write \(m=m_i\) and expand
\(B_i^h\) and \(B_{i+1}^h\) symmetrically about \(m\).  Using
\(B'=-\tau N\), the Frenet equations, \eqref{eq:chord-expansion},
and \eqref{eq:CB-explicit}, one obtains
\begin{align}
-\bigl\langle B_{i+1}^h,N_i^h\bigr\rangle
&=
h\tau(m)+h^3E(m)+O(h^4),
\label{eq:sinphi-expanded}\\
\bigl\langle B_{i+1}^h,B_i^h\bigr\rangle
&=
1-\frac{h^2}{2}\tau(m)^2+O(h^4),
\label{eq:cosphi-expanded}
\end{align}
where
\[
E=
\frac{\kappa^2\tau}{8}-\frac{\tau^3}{6}
+\frac{\tau''}{8}
+\frac{\kappa'\tau'}{6\kappa}
+\frac{\kappa''\tau}{6\kappa}
-\frac{(\kappa')^2\tau}{6\kappa^2}.
\]
Since our sign convention is
\[
\sin\phi_i^h=-\langle B_{i+1}^h,N_i^h\rangle,
\qquad
\cos\phi_i^h=\langle B_{i+1}^h,B_i^h\rangle,
\]
the \(\operatorname{atan2}\) expansion gives
\begin{equation}
\phi_i^h
=
h\tau(m_i)+h^3A_\phi(m_i)+O(h^4),
\label{eq:phi-refined}
\end{equation}
with
\[
A_\phi=
\frac{\kappa^2\tau}{8}
+\frac{\tau''}{8}
+\frac{\kappa'\tau'}{6\kappa}
+\frac{\kappa''\tau}{6\kappa}
-\frac{(\kappa')^2\tau}{6\kappa^2}.
\]
In particular,
\begin{equation}
\phi_i^h=h\tau(m_i)+O(h^3).
\label{eq:phi-asymptotic}
\end{equation}
The estimates \eqref{eq:x-asymptotic} and \eqref{eq:phi-asymptotic}, together with $\kappa\ge\kappa_0>0$ and $|\tau|\ge\tau_0>0$, imply after decreasing $h_0$ if necessary that $0<\theta_i^h<\pi$ and $\phi_i^h\notin\{0,\pi\}$ uniformly on the stated compact cells.

Consequently,
\[
\sin\phi_i^h=h\tau(m_i)+O(h^3),
\qquad
\cos\phi_i^h=1+O(h^2).
\]
Together with the refined $x_i^h$ expansion above and \eqref{eq:phi-refined}, symmetric midpoint expansion gives smooth functions $K$ and $L$ such that
\[
x_i^h+x_{i+1}^h\cos\phi_i^h
=
h\kappa(m_i)+h^3K(m_i)+O(h^4),
\]
and
\[
\sin\phi_i^h
=
h\tau(m_i)+h^3L(m_i)+O(h^4).
\]
Since \(|\tau|\ge\tau_0>0\), division yields
\begin{equation}
q_i^h:=
\frac{x_i^h+x_{i+1}^h\cos\phi_i^h}{\sin\phi_i^h}
=
\rho(m_i)+h^2S(m_i)+O(h^3),
\qquad
\rho=\frac{\kappa}{\tau},
\label{eq:q-refined}
\end{equation}
where
\[
S=\frac{K}{\tau}-\frac{\kappa L}{\tau^2}
\]
is smooth. In particular, $q_i^h=(\kappa/\tau)(m_i)+O(h^2)$, proving the stated second-order quotient consistency.

Since $m_{i+1}=m_i+h$,
\[
D_i^h=q_{i+1}^h-q_i^h
=\rho(m_i+h)-\rho(m_i)+O(h^2),
\]
and Taylor expansion yields
\[
\frac{D_i^h}{h}=\rho'(m_i)+O(h).
\]
If the underlying curve is a smooth general helix, then $\rho$ is constant and the refined expansion \eqref{eq:q-refined} gives
\[
D_i^h
=h^2\bigl(S(m_i+h)-S(m_i)\bigr)+O(h^3)
=O(h^3).
\]
Hence every smooth general helix satisfying the hypotheses has $D_i^h=O(h^3)$.
\end{proof}

\end{document}
